% Part: normal-modal-logic
% Chapter: tableaux
% Section: introduction

\documentclass[../../../include/open-logic-section]{subfiles}

\begin{document}

\olfileid{nml}{tab}{int}

\olsection{مقدمة}

ال!!^{tableau}s أشجار معينة (تتفرع إلى أسفل) من ال!!{signed
  formula}s؛ أي من أزواج يتكون كل منها من علامة لقيمة صدق ($\True$ أو
$\False$) ومن !!a{sentence}:
\[
\sFmla{\True}{!A} \text{ أو } \sFmla{\False}{!A}.
\]
تبدأ !!^a{tableau} بعدد من \emph{الافتراضات}. وتُولد كل
!!{signed formula} لاحقة بتطبيق إحدى قواعد الاستدلال. تضيف بعض قواعد
الاستدلال !!{signed formula}s واحدة أو أكثر إلى طرف من أطراف الشجرة، وتضيف
قواعد أخرى طرفين جديدين، فتنتج شعبتان. وتنتج القواعد !!{signed formula}s
تكون فيها ال!!{formula} أقل تعقيدًا من الصيغة في
ال!!{signed formula} التي طُبقت عليها القاعدة. وعندما تحوي شعبة كلًا من
$\sFmla{\True}{!A}$ و$\sFmla{\False}{!A}$، نقول إن الشعبة
\emph{مغلقة}. وإذا كانت كل شعبة في !!a{tableau} مغلقة، كانت
ال!!{tableau} كلها مغلقة. وتشكل ال!!{tableau} المغلقة !!a{derivation}
تبين أن مجموعة ال!!{signed formula}s التي بدأنا بها ال!!{tableau} غير
قابلة للإرضاء. ويمكن استعمال ذلك لتعريف علاقة~$\Proves$:
$\Gamma \Proves !A$ إذا وفقط إذا وجدت مجموعة منتهية
$\Gamma_0=\{!B_1,\dots,!B_n\}\subseteq\Gamma$ توجد لها !!{tableau}
مغلقة للافتراضات:
\[
\{\sFmla{\False}{!A}, \sFmla{\True}{!B_1}, \dots, \sFmla{\True}{!B_n}\}.
\]

في المنطق الجهوي، يلزم أن نوسع مفهوم !!{signed formula} ونضيف قواعد
تتناول~\iftag{prvBox}{$\Box$\iftag{prvDiamond}{ و
    $\Diamond$}{}}{$\Diamond$}. فإلى جانب العلامة ($\True$ أو
$\False$)، تحمل ال!!{formula}s في ال!!{tableau}s الجهوية أيضًا
\emph{بادئات}~$\sigma$. وهذه البادئات متتاليات غير فارغة من الأعداد
الصحيحة الموجبة؛ أي~$\sigma \in (\PosInt)^*\setminus\{\emptyseq\}$.
وعندما نكتب هذه البادئات نحذف الأقواس المحيطة~$\tuple{\ }$ ونفصل بين
ال!!{element}s المفردة بنقاط~$.$ بدلًا من الفواصل~$,$. وإذا كانت
$\sigma$ بادئة، فإن~$\sigma.n$ هي~$\sigma\concat\tuple{n}$؛ فمثلًا،
إذا كانت~$\sigma=1.2.1$، كانت~$\sigma.3$ هي~$1.2.1.3$. ومن ثم، مثلًا:
\[
\sFmla{\True}{\Box !A \lif !A}[1.2]
\]
هي \emph{!!{signed formula} ذات بادئة} (أو اختصارًا \emph{!!{formula}
  ذات بادئة}).

والحدس هو أن البادئة تسمي عالمًا في نموذج قد يرضي ال!!{formula}s الموجودة
على شعبة من !!a{tableau}؛ وإذا كانت~$\sigma$ تسمي عالمًا ما، فإن
$\sigma.n$ تسمي عالمًا يمكن النفاذ إليه من العالم الذي تسميه~$\sigma$.

\end{document}
